\section{The fundamental inequality of convexity}

Let X be a set; in the vector space \(R^{X}\) of all finite numerical functions \(^{1}\) defined on X, let P be the set of all positive real-valued functions on X. On the other hand, let M be a numerical function, \(^{2}\) finite or not, with values \(\geqslant0\) , defined on P, such that:

\(1^{\circ} M(0) = 0\) , and \(M\) is positively homogeneous, that is, \(M(\lambda f) = \lambda M(f)\) for every real number \(\lambda > 0\) .

\(2^{\circ} M\) is increasing in P, in other words the relation \(f \leqslant g\) implies \(M(f) \leqslant M(g)\) .

\(3^{\circ} M\) is convex in P, in other words (TVS, II, §2, No.~8) satisfies the relation \(M(f + g) \leqslant M(f) + M(g)\) .

Example. --- Suppose X is a finite set, for example the interval \([1,n]\) of N; denoting by \(x_{i}\) ( \(1 \leqslant i \leqslant n\) ) the coordinates of a vector \(x \in R^{n}\) , the functions \(M_{1}(\mathbf{x}) = \sum_{i=1}^{n} x_{i}\) and \(M_{\infty}(\mathbf{x}) = \sup_{1 \leqslant i \leqslant n} x_{i}\) satisfy the preceding conditions in the set P of vectors x with coordinates \(\geqslant 0\) .

Remark. --- Let S be a pointed convex cone contained in P (that is, a set such that \(S + S \subset S\) and \(\lambda S \subset S\) for \(\lambda \geqslant 0\) ; cf.~TVS, II, §2, No.~4); let M be a numerical function (finite or not) with values \(\geqslant 0\) , defined on S and satisfying in S the above conditions \(1^{\circ}\) , \(2^{\circ}\) and \(3^{\circ}\) . Then M can be extended to the entire set P, in such a way that the extended function (which we denote again by M) satisfies the same conditions: it suffices, for every function \(f \in P\) , to set \(M(f) = +\infty\) if there does not exist any function \(g \in S\) such that \(f \leqslant g\) , and \(M(f) = \inf_{g \in S, f \leqslant g} M(g)\) in

the contrary case. This procedure will be applied in Ch. IV, §1 to define the upper integral of a positive function.

PROPOSITION 1. --- Let \(\varphi(t_{1}, t_{2}, \ldots, t_{n})\) be a finite numerical function, defined and continuous for \(t_{i} \geqslant 0\) ( \(1 \leqslant i \leqslant n\) ), such that:

\(1^{\circ}\) the relations \(t_i > 0\) \((1\leqslant i\leqslant n)\) imply \(\varphi (t_1,t_2,\dots ,t_n) > 0\)

\(2^{\circ}\) the function \(\varphi\) is positively homogeneous;

\(3^{\circ}\) the set \(\mathbf{K} \subset \mathbf{R}^{n}\) defined by the relations \(t_{i} \geqslant 0\) ( \(1 \leqslant i \leqslant n\) ) \(\varphi(t_{1}, t_{2}, \ldots, t_{n}) \geqslant 1\) is convex.

Under these conditions, if \(f_1, f_2, \ldots, f_n\) are \(n\) finite functions \(\geqslant 0\) defined on \(X\) , such that \(M(f_i) < +\infty\) for \(1 \leqslant i \leqslant n\) , then

\[
M \left(\varphi \left(f _ {1}, f _ {2}, \dots , f _ {n}\right)\right) \leqslant \varphi \left(M \left(f _ {1}\right), M \left(f _ {2}\right), \dots , M \left(f _ {n}\right)\right).\tag{1}
\]

One knows, by the Hahn-Banach theorem (TVS, II, §5), that K is the intersection of the n half-spaces \(t_{i} \geqslant 0\) ( \(1 \leqslant i \leqslant n\) ) and a family of closed half-spaces \((\mathrm{U}_{\iota})_{\iota \in \mathrm{I}}\) , \(U_{\iota}\) being defined by a relation of the form

\[
\alpha_ {\iota 1} t _ {1} + \alpha_ {\iota 2} t _ {2} + \dots + \alpha_ {\iota n} t _ {n} - \beta_ {\iota} \geqslant 0,\tag{2}
\]

where the \(\alpha_{\iota k}\) are not all zero. By hypothesis, if \(\mathbf{t}=(t_{i})\) is such that \(t_{i}>0\) for \(1\leqslant i\leqslant n\) , then \(\varphi(t_{1},\ldots,t_{n})>0\) , therefore there exists a \(\lambda_{0}>0\) such that the relation \(\lambda\geqslant\lambda_{0}\) implies \(\lambda t\in K\) ; this shows that, for each \(\iota\in I\) , the relations \(t_{i}\geqslant0\) ( \(1\leqslant i\leqslant n\) ) imply \(\alpha_{\iota1}t_{1}+\cdots+\alpha_{\iota n}t_{n}\geqslant0\) and therefore \(\alpha_{\iota k}\geqslant0\) for \(1\leqslant k\leqslant n\) ; it is then clear that K is also the intersection of the half-spaces \(t_{i}\geqslant0\) ( \(1\leqslant i\leqslant n\) ) and the \(U_{\iota}\) such that \(\beta_{\iota}\geqslant0\) ; moreover, since the origin does not belong to K, there exists at least one index \(\iota\) such that \(\beta_{\iota}>0\) .

Now let C be the convex cone in \(R^{n+1}\) defined by the relations \(t_{i} \geqslant 0\) ( \(1 \leqslant i \leqslant n + 1\) ), \(t_{n+1} \leqslant \varphi(t_{1}, t_{2}, \ldots, t_{n})\) (the closure of the convex cone generated in \(R^{n+1}\) by the convex set \(K \times \{1\}\) ); it is immediate that C is also defined by the relations \(t_{i} \geqslant 0\) ( \(1 \leqslant i \leqslant n + 1\) ) and

\[
\beta_ {\iota} t _ {n + 1} \leqslant \alpha_ {\iota 1} t _ {1} + \dots + \alpha_ {\iota n} t _ {n} \quad (\iota \in \mathrm{I}, \beta_ {\iota} \geqslant 0).\tag{3}
\]

For every \(x \in X\) , we therefore have

\[
\beta_ {\iota} \varphi (f _ {1} (x), \dots , f _ {n} (x)) \leqslant \alpha_ {\iota 1} f _ {1} (x) + \dots + \alpha_ {\iota n} f _ {n} (x)\tag{4}
\]

for all \(\iota \in I\) . For every index \(\iota\) such that \(\beta_{\iota} > 0\) , it follows from (4) and the hypotheses on \(M\) that \(M(\varphi(f_1, f_2, \ldots, f_n))\) is finite and

\[
\beta_ {\iota} M \left(\varphi \left(f _ {1}, f _ {2}, \dots , f _ {n}\right)\right) \leqslant \alpha_ {\iota 1} M \left(f _ {1}\right) + \alpha_ {\iota 2} M \left(f _ {2}\right) + \dots + \alpha_ {\iota n} M \left(f _ {n}\right),
\]

and this relation is also verified in an obvious manner if \(\beta_{\iota} = 0\) . We thus

see that the point with coordinates

\[
M (f _ {1}), M (f _ {2}), \dots , M (f _ {n}), M (\varphi (f _ {1}, f _ {2}, \dots , f _ {n}))
\]

belongs to C, which proves the proposition.

